\section{Canonical comparison capacity at the original roof}
\label{sec:v108-pool}
The source and comparison trajectories need not be matched on an
entire angular block. We instead compare them at the roof at which
the source is evaluated. This distinction permits a canonical
construction, with physical overlap one, on the full comparison
space supplied by the inherited regular-segment hypotheses.

Fix the collision count, radius, exact labels and $\varepsilon>0$.
All integrals below are on the original disjoint physical pairing
space $\Theta_{\rm phys}$. We use its Borel segment endpoints and
the finite traces in \eqref{eq:v107-envelope}, together with
Lemma~\ref{lem:v107-endpoint}. These are precisely the inherited
regular-segment inputs, not new assertions about unverified singular
extensions. If those inputs have been verified only on a specified
Borel comparison subspace, every comparison integral in this section
is restricted to that subspace; the source outside it is retained.
Write $\ell=b-a$ and fix $0<r<\Delta$ before taking a collision limit.

\begin{definition}[The complete admissible pool]\label{def:v108-pool}
At the original roof $t$, put
\begin{align}
 q_+(t,\theta)&=\one_{\{0<\ell(\theta)<\Delta,\ a(\theta)<t,
                                      \ t+r<b(\theta)\}},\notag\\
 q_-(t,\theta)&=(1-q_+(t,\theta))
        \one_{\{0<\ell(\theta)<\Delta,\ a(\theta)<t-r,
                                      \ t<b(\theta)\}},\notag\\
 Q_r(t,\theta)&=q_+(t,\theta)+q_-(t,\theta).
                                                    \label{eq:v108-pool}
\end{align}
Thus a sheet is counted once even when both orientations are
available. Define
\begin{equation}\label{eq:v108-PDE}
 P_r(t)=\int Q_r p_\vee\,d\theta,\qquad
 D_r(t)=\int Q_r c^\varepsilon\,d\theta,\qquad
 E_r(t)=\int Q_r e_\partial\,d\theta.
\end{equation}
In particular $P_r=D_r+E_r$. The pool contains every admissible
physical segment containing one of the two closed roof windows
$[t,t+r]$ and $[t-r,t]$ in its interior. It is not a selected
collection of angular source blocks.
\end{definition}

\begin{lemma}[Measurability and one-copy localization]
\label{lem:v108-pool}
The functions in \eqref{eq:v108-pool} are jointly Borel and satisfy
$0\le Q_r\le1$. Their integrals have common measurable roof versions
at the fixed collision count, radius and exact labels. On
$\{q_+>0\}$ the segment contains $[t,t+r]$; on $\{q_->0\}$ it
contains $[t-r,t]$. On $\{Q_r>0\}$ its entire interior lies in
$(t-\Delta,t+\Delta)$.
\end{lemma}
\begin{proof}
The indicator conditions are Borel inequalities between the inherited
endpoint functions and $t$. The factor $1-q_+$ disjointizes the two
orientations. Integration of a nonnegative jointly Borel function
against the fixed sigma-finite angular measure gives a measurable
function of $t$. The asserted window inclusions follow from the
strict inequalities. Finally $a<t<b$ and $b-a<\Delta$ imply
$t-\Delta<a<b<t+\Delta$. All statements concern the original
physical measure, so a disjoint refinement of its charts changes
none of the integrals. We use the inherited common coarea versions;
a null set uniform over all radii is neither needed nor asserted.
\end{proof}

\begin{theorem}[Canonical finite-count envelope bound]
\label{thm:v108-pool}
At finite count, on the original roof versions,
\begin{equation}\label{eq:v108-capacity-bound}
 P_r(t)\le \frac1r\int_{t-r}^{t+r}\mathsf T_\Delta^\varepsilon(u)\,du
       +C(1+\varepsilon^{-1})\int_{t-\Delta}^{t+\Delta}v^0(u)\,du.
\end{equation}
Consequently, for any previously fixed $K>0$,
\begin{equation}\label{eq:v108-capacity-height}
 \mathcal H_M(KP_r)\le
 CK\{\varepsilon^{1/16}+(1+\varepsilon^{-1})\Delta\}.
\end{equation}
No source-block balance, angular continuation, monotonicity of either
endpoint roof, or independently supplied overlap bound is required.
\end{theorem}
\begin{proof}
For the positive orientation, Tonelli's theorem gives
\[
 r\int q_+(t,\theta)c^\varepsilon(\theta)\,d\theta
 =\int_t^{t+r}\int q_+(t,\theta)c^\varepsilon(\theta)
                  \one_{\{a(\theta)<u<b(\theta)\}}\,d\theta\,du.
\]
The analogous identity holds on $(t-r,t)$ for $q_-$. Dropping the
comparison weights on their respective half-windows proves
\[
 D_r(t)\le\frac1r\int_{t-r}^{t+r}
                         \mathsf T_\Delta^\varepsilon(u)\,du.
\]
The weights depend on $t$, but $t$ is fixed throughout these
integrals; they are never evaluated at $u$ or differentiated.
Lemma~\ref{lem:v107-endpoint} and
Lemma~\ref{lem:v108-pool} give
\begin{align*}
 E_r(t)&\le C(1+\varepsilon^{-1})
   \int Q_r(t,\theta)\int_{a(\theta)}^{b(\theta)}
                 A(u,\theta)I_m(u,\theta)/z(u,\theta)\,du\,d\theta\\
 &\le C(1+\varepsilon^{-1})
                         \int_{t-\Delta}^{t+\Delta}v^0(u)\,du.
\end{align*}
The last step uses the original one-copy physical summation and
$Q_r\le1$, not disjointness of auxiliary angular blocks. Adding the
two inequalities proves \eqref{eq:v108-capacity-bound}.
All sums are nonnegative and are completed before any norm or
collision limit. The inherited fixed-count local integrability makes
the displayed quantities finite for almost every roof. Finally use
\eqref{eq:v104-paired-average} at $r$ and
\eqref{eq:v104-uncut-input} at $\Delta$. The factors are respectively
$2K$ and $2CK(1+\varepsilon^{-1})\Delta$. This proves
\eqref{eq:v108-capacity-height} in the prescribed order of limits.
\end{proof}

Let $0\le H\le1$ be an original Borel selector and put
\[
 U(t)=\int c^\varepsilon(\theta)\one_{\{a(\theta)<t<b(\theta)\}}
                         H(\Psi(\theta,t))\,d\theta.
\]
Assume the inherited local finiteness of this positive source. Set
\begin{equation}\label{eq:v108-allocation}
 \begin{aligned}
 V_{r,K}(t)&=\min\{U(t),KP_r(t)\},&
 R_{r,K}(t)&=(U(t)-KP_r(t))_+,\\
 \beta_{r,K}(t)&=\begin{cases}
    \min\{1,KP_r(t)/U(t)\},&U(t)>0,\\
    0,&U(t)=0.
   \end{cases}
 \end{aligned}
\end{equation}
Definitions may be completed arbitrarily on the common roof-null
set where finiteness fails.

\begin{corollary}[An exact original-source allocation]
\label{cor:v108-allocation}
The multiplier $\beta_{r,K}$ is measurable and takes values in
$[0,1]$. At finite count,
\begin{equation}\label{eq:v108-source-split}
 U=V_{r,K}+R_{r,K},\qquad
 V_{r,K}(t)=\int c^\varepsilon\one_{\{a<t<b\}}
                    H(\Psi(\theta,t))\beta_{r,K}(t)\,d\theta.
\end{equation}
Replacing $\beta_{r,K}$ by $1-\beta_{r,K}$ gives the corresponding
integral formula for $R_{r,K}$. Moreover $V_{r,K}$ satisfies the
right sides of \eqref{eq:v108-capacity-bound} and
\eqref{eq:v108-capacity-height}, multiplied by $K$ in the former.
\end{corollary}
\begin{proof}
Measurability follows from Lemma~\ref{lem:v108-pool}. The identities
follow pointwise from the minimum and positive-part formulas.
Since $V_{r,K}\le KP_r$, Theorem~\ref{thm:v108-pool} applies.
The endpoint maximum supplies only comparison capacity. The measure
on the source side still has its original minimum trace and its
original target selector; neither has been moved to another roof.
\end{proof}
